\chapter{复流形与Kahler几何}
\label{ch:17-complex}
\ChapterPlan{周第7章（本地页码待核对）；Moroianu讲义第1--7节补充}{正文证明初稿；8题详解及3幅原创图，尚未独立数学审阅}

前一章中的复向量丛只要求纤维具有复线性运算。本章进一步要求底空间的坐标变换全纯，因而切空间上的乘$i$能与微分、度量及联络相互制约。我们从实流形出发，依次讨论几乎复结构、可积性、Hermite度量和Kahler条件；不能将这四层条件混成一个定义。

本章流形仍有限维、Hausdorff、第二可数、无边界，复维记为$m\geq1$，实维为$2m$。实曲率沿用$R(X,Y)Z=\nabla_X\nabla_YZ-\nabla_Y\nabla_XZ-\nabla_{[X,Y]}Z$。所用单变量Cauchy积分公式是明确的复分析先修，不在本书重新证明；Newlander--Nirenberg的充分性和一般局部势存在另作外引，不让第14章实Hodge理论承担尚未建立的复分析结论。

\Needspace{16\baselineskip}
\section{复结构与全纯坐标}
\subsection{从全纯图册到几乎复张量}
\begin{definition}\label{def:complex-atlas}
实$2m$维流形的复图册由取值于$\mathbb C^m$开集的坐标卡组成，重叠处换元为全纯双射且逆全纯。两个图册的并仍满足此条件时视为同一复结构，装备复结构者称复流形（complex manifold）。几乎复结构（almost complex structure）则仅为光滑实丛自同态$J:TM\to TM$，满足$J^2=-I$。
\end{definition}
这里全纯可理解为在每点具有收敛复幂级数表示。对光滑映射$f:U\subset\mathbb C^m\to\mathbb C^r$，等价判据是$\partial f/\partial\bar z^j=0$：幂级数逐项求导给必要性；反向固定除一个以外的变量，用单变量Cauchy公式，依次在紧含于$U$的多圆盘上迭代，得到
\[
 f(z)=\frac1{(2\pi i)^m}\int\cdots\int
 \frac{f(\zeta)\,\dd\zeta^1\cdots\dd\zeta^m}
 {\prod_j(\zeta^j-z^j)}.
\]
把各分母按几何级数展开，在更小闭多圆盘上一致收敛，故积分可逐项进行并给出幂级数。这一步使用单变量复分析及紧集上一致收敛的积分换序，没有证明一般几乎复结构存在这类坐标。

写$z^j=x^j+iy^j$，规定
\[
 J\partial_{x^j}=\partial_{y^j},\quad
 J\partial_{y^j}=-\partial_{x^j},\qquad
 \partial_j=\tfrac12(\partial_{x^j}-i\partial_{y^j}),\quad
 \partial_{\bar j}=\tfrac12(\partial_{x^j}+i\partial_{y^j}).
\]
全纯换元的实微分与乘$i$交换，所以局部$J$拼成全局张量，且$J^2=-I$。反过来，若$J$确能在一族实坐标中写成上述形式，则两张这种坐标间的微分与乘$i$交换；这就是Cauchy--Riemann方程，故换元全纯。这只能证明\emph{已有适配坐标之间}的相容性，不能从$J^2=-I$直接造出坐标。

\Needspace{14\baselineskip}
\begin{proposition}\label{prop:complex-map-inverse}
复流形间的光滑映射$f$全纯，当且仅当$\dd f\circ J_M=J_N\circ\dd f$。若复维相同且$\dd f_p$可逆，则$f$在$p$附近为双全纯映射。
\end{proposition}
\begin{proof}
坐标下微分与乘$i$交换等价于各分量的$\partial_{\bar j}f=0$，由上述判据得第一条。实逆映射定理给光滑局部逆；等式左右乘逆微分说明逆也与$J$交换，故逆全纯。这里没有把任意实局部微分同胚误认为双全纯。
\end{proof}
例如共轭映射$z\mapsto\bar z$为实微分同胚，但其微分反交换$J$，不全纯。复环面$\mathbb C^m/\Lambda$中$\Lambda$为实满秩格：足够小的开球与非零格平移不交，投影给局部坐标；重叠换元为平移，故全纯。格离散和有界基本平行多面体分别给Hausdorff性与紧致性。

\subsection{复定向及线性类型分解}
复矩阵$A+iB$在实坐标中的矩阵为$\left(\begin{smallmatrix}A&-B\\B&A\end{smallmatrix}\right)$，复化后相似于$\operatorname{diag}(A+iB,A-iB)$。因此实行列式为$|\det_{\mathbb C}(A+iB)|^2>0$，全纯坐标自然定向。一般几乎复流形也可在每点附近选复基$v_1,\ldots,v_m$，以$(v_1,Jv_1,\ldots,v_m,Jv_m)$定向；复换基仍有正实行列式，故定向拼接。

在$T_{\mathbb C}M=TM\otimes_{\mathbb R}\mathbb C$上复线性延拓$J$。多项式$u^2+1$的根互异，因而
\begin{equation}\label{eq:complex-projections}
 T_{\mathbb C}M=T^{1,0}M\oplus T^{0,1}M,\qquad
 \pi^{1,0}=\tfrac12(I-iJ),\quad\pi^{0,1}=\tfrac12(I+iJ).
\end{equation}
两空间分别为$+i,-i$特征丛，互为复共轭，秩均为$m$。在全纯坐标中分别由$\partial_j,\partial_{\bar j}$张成。余切丛的$(1,0)$部分定义为$T^{0,1}$的零化子，局部基为$\dd z^j$，另一部分为$\dd\bar z^j$。外幂分解为$\Lambda^kT^*_{\mathbb C}M=\bigoplus_{p+q=k}\Lambda^{p,q}$；基中先列$p$个$\dd z$再列$q$个$\dd\bar z$，不额外加入阶乘。

\begin{figure}[htbp]
\centering\DGDiagram{complex-types}

\caption{复化切丛的两个特征子丛。箭头表示代数投影与共轭；共轭为反线性，图不是把实切空间分成两个实子空间。}
\label{fig:complex-types}
\end{figure}

\subsection{Dolbeault算子及其成立条件}
在复流形上，按全纯坐标定义
\begin{align}
 \partial(f\,\dd z^I\wedge\dd\bar z^J)
 &=\sum_a\partial_a f\,\dd z^a\wedge\dd z^I\wedge\dd\bar z^J,\notag\\
 \bar\partial(f\,\dd z^I\wedge\dd\bar z^J)
 &=\sum_a\partial_{\bar a}f\,\dd\bar z^a\wedge\dd z^I\wedge\dd\bar z^J.
 \label{eq:dolbeault-local}
\end{align}
第二式如需重排为约定基，须移动$\dd\bar z^a$穿过全部$\dd z^I$，符号为$(-1)^{|I|}$。全纯换元不混合$\dd z$和$\dd\bar z$，故这两个算子就是$\dd$的$(p+1,q)$和$(p,q+1)$投影，与坐标无关。

\Needspace{14\baselineskip}
\begin{proposition}\label{prop:dolbeault-complex}
复流形上$\dd=\partial+\bar\partial$，且
\[
 \partial^2=0,\qquad\bar\partial^2=0,\qquad
 \partial\bar\partial+\bar\partial\partial=0.
\]
两算子都按总次数满足分次Leibniz规则，复共轭交换$\partial$和$\bar\partial$。
\end{proposition}
\begin{proof}
在坐标基上$\dd\dd z^j=\dd\dd\bar z^j=0$，所以只需对系数求导，得到分解。再把$\dd^2=0$分别投影到三种不同型别，得三条恒等式；或者直接用二阶偏导可交换和楔积反对称验证。分次乘法式与共轭关系也逐个坐标单项验证，因已知自然性而成为全局结论。
\end{proof}
例如$\partial\bar\partial f=\sum_{a,b}\partial_a\partial_{\bar b}f\,\dd z^a\wedge\dd\bar z^b$，而$\bar\partial\partial f$为其负号。定义
\[
 H^{p,q}_{\bar\partial}(M)=
 \frac{\ker(\bar\partial:\Omega^{p,q}\to\Omega^{p,q+1})}
 {\operatorname{im}(\bar\partial:\Omega^{p,q-1}\to\Omega^{p,q})},
\]
分母在$q=0$时为零。这里仅定义Dolbeault上同调，不预先宣称它与某个层上同调相同，亦不从第14章自动推出$H^k_{\mathrm{dR}}=\bigoplus_{p+q=k}H^{p,q}_{\bar\partial}$。

在仅有几乎复结构时，投影仍然存在，但$\dd$还可能含$(p+2,q-1)$和$(p-1,q+2)$分量。因为对一形式$\alpha$有$\dd\alpha(U,V)=U\alpha(V)-V\alpha(U)-\alpha([U,V])$，若$U,V\in T^{0,1}$而$[U,V]$出现$(1,0)$分量，$(1,0)$形式的外微分便出现$(0,2)$项。第17.4节将给出显式反例，故不能在几乎复情形不加条件地沿用本命题。

\subsection{全纯向量丛及其自然算子}
全纯向量丛是过渡矩阵全纯的复向量丛。取全纯局部标架$e$，对系数列$a$定义$\bar\partial_E(ea)=e\bar\partial a$。若$e'=eG$，全纯性$\bar\partial G=0$使定义相容，且$\bar\partial_E^2=0$。$T^{1,0}M$由全纯坐标标架及其Jacobi矩阵成为全纯丛，其对偶与外幂也一样；$K_M=\Lambda^{m,0}T^*M$称典范线丛。

反向，“光滑丛上的一个$(0,1)$联络满足平方零，便能找到全纯标架”是Koszul--Malgrange型可积性结论，本章不证明也不在构造中使用。下节Chern联络的起点始终是\emph{已经给定}的全纯丛，而不是仅给一族形式后默认为已存在全纯结构。


\Needspace{16\baselineskip}
\section{Hermite度量与Kahler形式}
\subsection{正定度量、基本形式及归一化}
\begin{definition}\label{def:hermitian-kahler}
几乎复流形上的Hermite度量为满足$g(JX,JY)=g(X,Y)$的黎曼度量。基本形式取
\[
 \omega(X,Y)=g(JX,Y).
\]
若$J$来自复图册且$\dd\omega=0$，称$(M,J,g)$为Kahler流形。
\end{definition}
由$J^2=-I$及度量相容可得$g(JX,Y)=-g(X,JY)$，所以$\omega$确为实二形式。若只有$J^2=-I$和$\dd\omega=0$，通常称几乎Kahler结构；不能在定义中略去可积性。

\begin{proposition}\label{prop:hermitian-exists}
每个几乎复流形都有Hermite度量；其基本形式非退化、型别为$(1,1)$。按复定向有$\omega^m/m!=\mathrm{vol}_g$。
\end{proposition}
\begin{proof}
由第11章取黎曼度量$g_0$，令$g(X,Y)=(g_0(X,Y)+g_0(JX,JY))/2$，正定且$J$不变。若$\omega(X,Y)=0$对所有$Y$成立，取$Y=JX$便得$g(X,X)=0$，所以$X=0$。$\omega(JX,JY)=\omega(X,Y)$，作用于两个同型特征向量时右式又乘$i^2=-1$或$(-i)^2=-1$，故同型部分为零，只余$(1,1)$。

选$e_1,Je_1,\ldots,e_m,Je_m$为实正交归一标架，记对偶为$u^j,v^j$。逐对代入得$\omega=\sum_j u^j\wedge v^j$；其$m$次楔积只有使用每对一次的$m!$项，符号均正，故给出体积公式。
\end{proof}
在全纯坐标中令$G=(g_{a\bar b})$，其中$g_{a\bar b}=g_{\mathbb C}(\partial_a,\partial_{\bar b})$，$g_{\mathbb C}$为$g$的\emph{复双线性}延拓。则
\begin{equation}\label{eq:hermitian-normalization}
 g=2\operatorname{Re}\sum_{a,b}g_{a\bar b}\,\dd z^a\otimes\dd\bar z^b,
 \qquad\omega=i\sum_{a,b}g_{a\bar b}\,\dd z^a\wedge\dd\bar z^b.
\end{equation}
Euclid度量$\sum_j((\dd x^j)^2+(\dd y^j)^2)$对应$G=I/2$和$\omega=\sum_j\dd x^j\wedge\dd y^j$，这个$1/2$以后不再随意改动。对$T^{1,0}M$取“第一变量共轭线性”的纤维Hermite内积时，其矩阵$H$为$H_{\bar a b}=g_{b\bar a}$，即$H=G^{\mathsf T}$；不能在矩阵联络公式中漏掉这次转置。

\subsection{闭性与Levi--Civita联络}
\begin{theorem}\label{thm:kahler-parallel}
在已经可积的Hermite流形上，以下条件等价：$\dd\omega=0$；
\[
 \partial_a g_{b\bar c}=\partial_b g_{a\bar c}
 \quad\text{及其复共轭式};
\]
Levi--Civita联络满足$\nabla J=0$。此时非零的同型联络系数为
\begin{equation}\label{eq:kahler-christoffel}
 \nabla_{\partial_a}\partial_b=\Gamma^c_{ab}\partial_c,
 \quad\Gamma^c_{ab}=\sum_d g^{\bar d c}\partial_a g_{b\bar d},
 \quad\nabla_{\partial_{\bar a}}\partial_b=0,
\end{equation}
以及它们的共轭；这里$\sum_dg_{b\bar d}g^{\bar d c}=\delta_b^c$。
\end{theorem}
\begin{proof}
展开$\partial\omega$，$\dd z^a\wedge\dd z^b\wedge\dd\bar z^c$的系数为$i(\partial_ag_{b\bar c}-\partial_bg_{a\bar c})$（对$a<b$求和）；$\bar\partial\omega$给共轭条件，且型别不同不能互相抵消。

在可交换的复坐标场上使用第11章Koszul公式。因两个同型向量的度量配对为零，得到
\begin{align*}
 2g_{\mathbb C}(\nabla_{\partial_a}\partial_b,\partial_{\bar c})
 &=\partial_ag_{b\bar c}+\partial_bg_{a\bar c},\\
 2g_{\mathbb C}(\nabla_{\partial_{\bar a}}\partial_b,\partial_{\bar c})
 &=\partial_{\bar a}g_{b\bar c}-\partial_{\bar c}g_{b\bar a},\\
 2g_{\mathbb C}(\nabla_{\partial_{\bar a}}\partial_b,\partial_c)
 &=\partial_bg_{c\bar a}-\partial_cg_{b\bar a}.
\end{align*}
另外$g_{\mathbb C}(\nabla_{\partial_a}\partial_b,\partial_c)=0$。代入闭性条件，混合导数同时与两型全部基向量配对为零，因此为零；第一式则化为\eqref{eq:kahler-christoffel}。联络保持两特征丛，等价于$\nabla J=0$。反向若$\nabla J=0$，结合$\nabla g=0$得$\nabla\omega=0$；无挠性使$\dd\omega$为$\nabla\omega$的交替和，所以闭。
\end{proof}
复维一时三形式恒零，因此任意Hermite度量均为Kahler。复维$m\geq2$时，$\widetilde g=e^{2u}g$有$\dd\widetilde\omega=2e^{2u}\dd u\wedge\omega$；在一点取上述正交基，映射$\alpha\mapsto\alpha\wedge\omega^{m-1}$为一次形式到$(2m-1)$次形式的同构。因此$\dd u\wedge\omega=0$迫使$\dd u=0$。例如$\mathbb C^2$上$e^{2x^1}g_{\mathrm{Euclid}}$是Hermite而非Kahler。

\subsection{势函数能给出什么}
若实光滑函数$\varphi$的复Hessian$(\partial_a\partial_{\bar b}\varphi)$正定，称其严格多重次调和函数（strictly plurisubharmonic）。取
\[
 \omega_\varphi=i\partial\bar\partial\varphi,
 \qquad g_{a\bar b}=\partial_a\partial_{\bar b}\varphi,
\]
则$\overline{\omega_\varphi}=\omega_\varphi$，且由两Dolbeault平方零及反交换可得$\dd\omega_\varphi=0$；正定性给出Kahler度量。$\varphi=|z|^2/2$恢复Euclid模型。若$f$为全纯函数，$\varphi+\operatorname{Re}f$给同一形式，因为$\partial\bar\partial f=\partial\bar\partial\bar f=0$。

\textbf{外引边界。}一般Kahler形式在每点附近都能写为$i\partial\bar\partial\varphi$，这是局部$\partial\bar\partial$引理的结论。本章仅证明上述“给定势产生度量”的方向，并显式构造下面的势；一般逆向存在性依赖尚未补证的局部Dolbeault求解，不以实Poincare引理一行代替。更不声称存在全局实势：紧致情形会受到Stokes障碍。

\subsection{复射影空间与Fubini--Study形式}
记$\mathbb{CP}^m$为$\mathbb C^{m+1}$中复直线的集合，取非零齐次坐标$[Z_0:\cdots:Z_m]$。在$U_j=\{Z_j\ne0\}$中用$z^a=Z_a/Z_j$（省略$j$）坐标，逆为在第$j$位填1。重叠换元为无零分母的有理全纯式，所以给复图册。拓扑取非零向量商拓扑；映射$[Z]\mapsto ZZ^*/|Z|^2$把它连续单射到Hermite矩阵空间，坐标逆在每个$U_j$上连续。这给出Hausdorff性；有限张可数基坐标卡给第二可数性。单位球面满射到商空间，故该流形紧致。$m=1$时正是第16章构造的底空间。

在第$j$张卡上令$\varphi_j=\log(\sum_a|Z_a|^2/|Z_j|^2)$。重叠处
\[
 \varphi_k-\varphi_j=-\log|Z_k/Z_j|^2.
\]
对任意无零全纯$f$，$\partial\log|f|^2=f^{-1}\partial f$，故$\bar\partial\partial\log|f|^2=0$，无须选择全局对数分支。于是$i\partial\bar\partial\varphi_j$在重叠处相等，拼成全局$\omega_{\mathrm{FS}}$。

\Needspace{18\baselineskip}
\begin{theorem}\label{thm:fubini-study}
上述形式为Kahler形式。在$U_0\simeq\mathbb C^m$上，记$r^2=\sum_a|z^a|^2$，其系数及行列式为
\begin{equation}\label{eq:fs-matrix}
 g_{a\bar b}=\frac{(1+r^2)\delta_{ab}-\bar z^a z^b}{(1+r^2)^2},
 \qquad\det G=(1+r^2)^{-(m+1)}.
\end{equation}
本书的Fubini--Study度量取$g_{\mathrm{FS}}=2\operatorname{Re}(G_{a\bar b}\dd z^a\otimes\dd\bar z^b)$。
\end{theorem}
\begin{proof}
直接对$\log(1+r^2)$作两次Wirtinger求导即得系数。对$v\ne0$，Cauchy--Schwarz不等式给
\[
 \sum_{a,b}g_{a\bar b}v^a\bar v^b
 =\frac{(1+r^2)|v|^2-|\sum_a\bar z^a v^a|^2}{(1+r^2)^2}
 \geq\frac{|v|^2}{(1+r^2)^2}>0.
\]
因此正定。闭性由势函数已经证明，且局部形式已相容。行列式可在$z\ne0$时沿$\bar z$及其正交补分解矩阵：两个特征值分别为$(1+r^2)^{-2}$和$(1+r^2)^{-1}$，后者重数$m-1$；相乘给公式，在$z=0$处由连续性也成立。
\end{proof}

\begin{figure}[htbp]
\centering\DGDiagram{projective-potentials}

\caption{局部势在两张射影卡上不同，但其$i\partial\bar\partial$相同。图表示重叠区上的公式关系；没有全局实势的断言，亦不是射影空间的等距平面图。}
\label{fig:projective-potentials}
\end{figure}

在$\mathbb{CP}^1$上，$w=x+iy$给
\[
 g_{\mathrm{FS}}=\frac{2(\dd x^2+\dd y^2)}{(1+|w|^2)^2},\qquad
 \omega_{\mathrm{FS}}=\frac{2\dd x\wedge\dd y}{(1+|w|^2)^2},
 \quad\int_{\mathbb{CP}^1}\omega_{\mathrm{FS}}=2\pi.
\]
第16章单位球面识别的拉回度量为$4(\dd x^2+\dd y^2)/(1+|w|^2)^2$，故这里半径为$1/\sqrt2$、Gauss曲率为2；Hopf曲率恰为$F=i\omega_{\mathrm{FS}}$。这组常数与有些书采用$i\partial\bar\partial\log(1+r^2)/2$的约定不同，比较曲率时必须同时比较尺度。

全纯浸入$f:N\to M$拉回Kahler形式仍闭，且$f^*\omega(X,JX)=g(\dd fX,\dd fX)>0$，所以诱导度量为Kahler。一般全纯映射有临界点时，拉回仅半正定，不一定是度量。两个Kahler流形的乘积取两形式拉回之和，闭性及分块正定性直接给Kahler结构；例如标准复环面及$\mathbb{CP}^1\times\mathbb{CP}^1$均给出紧例子。


\Needspace{16\baselineskip}
\section{复流形上的联络与曲率}
\subsection{Chern联络的构造及唯一性}
设$E\to M$为\emph{已经给定}的全纯向量丛，纤维Hermite内积$h$取第一变量共轭线性、第二变量线性。全纯标架$e=(e_1,\ldots,e_r)$中
\[
 h(ea,eb)=\bar a^{\mathsf T}Hb,\qquad H=\bar H^{\mathsf T}>0,
 \qquad\nabla(ea)=e(\dd a+Aa).
\]
对矩阵值形式记$A^*=\bar A^{\mathsf T}$，共轭同时交换$(1,0)$与$(0,1)$型。

\begin{theorem}\label{thm:chern-connection}
存在唯一保持$h$且$(0,1)$部分等于$\bar\partial_E$的复线性联络，称Chern联络。其全纯标架表达式、曲率为
\begin{equation}\label{eq:chern-matrix}
 A=H^{-1}\partial H,\qquad
 F=\bar\partial(H^{-1}\partial H),\qquad F^{2,0}=F^{0,2}=0.
\end{equation}
若$e'=eG$为全纯换标架，则$H'=G^*HG$，且$A'=G^{-1}AG+G^{-1}\dd G$。
\end{theorem}
\begin{proof}
度量相容在实向量输入上等价于$\dd H=A^*H+HA$，由对$h(ea,eb)$求导并消去$\dd a,\dd b$项得到。条件$\nabla^{0,1}=\bar\partial_E$要求$A^{0,1}=0$，所以该等式的$(1,0)$部分为$\partial H=HA$，唯一迫使$A=H^{-1}\partial H$。反向取此$A$，其共轭转置满足$A^*H=\bar\partial H$，两型合起来给度量相容。

换全纯标架时$\partial G^*=0$，代入$H'=G^*HG$并展开可得$A'=G^{-1}AG+G^{-1}\partial G$。又$\bar\partial G=0$，故这就是第10、16章换标架律，局部联络可拼接。曲率的$(0,2)$部分为零；利用$\partial H^{-1}=-H^{-1}(\partial H)H^{-1}$，
\[
 \partial A=-H^{-1}(\partial H)H^{-1}\wedge\partial H=-A\wedge A.
\]
故$(2,0)$部分也为零，只剩$F=\bar\partial A$，并按$F'=G^{-1}FG$变换。这里$\bar\partial\partial=-\partial\bar\partial$，曲率号不可换掉。
\end{proof}
全纯线丛时$A=\partial\log H$、$F=-\partial\bar\partial\log H$。例如平凡线丛配$H=e^{|z|^2}$，便有$F=-\dd z\wedge\dd\bar z$；丛平凡不意味着其给定度量的Chern联络平坦。若只要求度量相容而不固定$(0,1)$部分，则第16章表明还可加适当的斜Hermite一形式，唯一性会失去。

\subsection{与切丛的Levi--Civita联络比较}
对$E=T^{1,0}M$，按\eqref{eq:hermitian-normalization}使用$H=G^{\mathsf T}$，Chern系数为
\[
 A^c{}_b=\sum_{a,d}g^{\bar d c}\partial_a g_{b\bar d}\,\dd z^a.
\]
在Kahler情形，这恰与\eqref{eq:kahler-christoffel}一致，故Levi--Civita复化后在$T^{1,0}$上的限制为Chern联络。反向，若实Levi--Civita联络复化后保持$T^{1,0}$并与该Chern联络相同，由复共轭也保持$T^{0,1}$，所以$\nabla J=0$，进而Kahler。一般Hermite度量的Levi--Civita未必保持$T^{1,0}$，不能不加条件地把其“限制”叫作复丛联络。

把Chern联络及其共轭合为$T_{\mathbb C}M$的联络，则其$(1,0)$方向的挠率为
\[
 T^c{}_{ab}=\sum_dg^{\bar d c}
 (\partial_ag_{b\bar d}-\partial_bg_{a\bar d}).
\]
混合挠率为零，另一同型部分是共轭；因此无挠恰等价于Kahler闭性。这个比较解释了为何同一Hermite度量会自然产生两个不同联络。

\subsection{迹曲率、Ricci形式与尺度}
对Chern矩阵取迹，Jacobi行列式求导公式给$\operatorname{tr}A=\partial\log\det H$。因为$\det H>0$，这里用普通实对数，无须额外的线丛平凡化。于是
\begin{equation}\label{eq:chern-trace}
 \operatorname{tr}F=-\partial\bar\partial\log\det H.
\end{equation}
换标架的曲率共轭律使迹成为全局二形式，且度量相容使其在实向量输入上为纯虚，故$i\operatorname{tr}F$为实形式。两Hermite度量的行列式之比在换标架时不变，因此它们的迹曲率相差全局$-\partial\bar\partial$的实函数；这里只证明此微分形式结论，第18章将构造实de Rham示性类，一般整类的构造与比较另列待办。

\begin{theorem}\label{thm:kahler-ricci}
Kahler流形的Ricci形式取$\rho(X,Y)=\operatorname{Ric}(JX,Y)$，则
\begin{equation}\label{eq:kahler-ricci}
 \rho=i\operatorname{tr}F_{T^{1,0}}
 =-i\partial\bar\partial\log\det G,
 \qquad \operatorname{Ric}_{a\bar b}=-\partial_a\partial_{\bar b}\log\det G.
\end{equation}
特别地$\dd\rho=0$。本书的Fubini--Study度量满足$\rho=(m+1)\omega_{\mathrm{FS}}$，即实张量$\operatorname{Ric}=(m+1)g_{\mathrm{FS}}$。
\end{theorem}
\begin{proof}
在全纯坐标中，$\nabla_{\bar b}\partial_c=0$，所以
\[
 R(\partial_a,\partial_{\bar b})\partial_c
 =-\partial_{\bar b}\Gamma^d_{ac}\partial_d.
\]
Chern曲率为$(1,1)$型说明$R(\partial_a,\partial_b)=0$及其共轭。计算第11章的$\operatorname{Ric}(X,Y)=\operatorname{tr}(Z\mapsto R(Z,X)Y)$时，可在复化空间取迹，数值等于实迹的复双线性延拓。取$X=\partial_a,Y=\partial_{\bar b}$，在$T^{1,0}$上的对角块为零；在$T^{0,1}$上有
\[
 R(\partial_{\bar k},\partial_a)\partial_{\bar b}
 =-\partial_a\overline{\Gamma^l_{kb}}\,\partial_{\bar l}.
\]
挠率为零使$\Gamma^l_{kb}=\Gamma^l_{bk}$，故其迹为$-\partial_a\partial_{\bar b}\log\det G$，两个同型Ricci分量为零。$J\partial_a=i\partial_a$给出$\rho$的系数，正好是\eqref{eq:chern-trace}乘$i$。闭性由Dolbeault恒等式得到。最后将\eqref{eq:fs-matrix}的行列式代入，便得$\rho=(m+1)i\partial\bar\partial\log(1+r^2)$；由$g(X,Y)=\omega(X,JY)$还原实Ricci张量即可。
\end{proof}
将度量乘正常数$c$，$G$乘$c$，$\log\det G$增加常数$m\log c$，所以$\rho$与作为$(0,2)$张量的Ricci不变，而Einstein常数除以$c$。复维一时$\rho=K\omega$，故前面的$\mathbb{CP}^1$模型确有$K=2$。一般Hermite流形上$i\operatorname{tr}F$仍有定义，但不能未经Kahler假设便把它当作实Levi--Civita的Ricci形式。

\subsection{紧致性产生的上同调障碍}
\begin{proposition}\label{prop:kahler-powers}
紧致无边界Kahler流形的$[\omega]^k\in H^{2k}_{\mathrm{dR}}(M;\mathbb R)$在$1\leq k\leq m$时均非零。特别地其第二de Rham上同调非零，且不存在全局实势$\omega=i\partial\bar\partial\varphi$。
\end{proposition}
\begin{proof}
可逐个连通分支证明。复定向及体积公式给$\int_M\omega^m=m!\operatorname{Vol}_g(M)>0$。若$\omega^k=\dd\eta$，闭性使$\omega^m=\dd(\eta\wedge\omega^{m-k})$，与Stokes矛盾。若有全局实势，则$\omega=\dd\bigl(\tfrac i2(\bar\partial-\partial)\varphi\bigr)$；括号内为实一形式，又与$k=1$结论矛盾。非紧空间上的Euclid全局势不矛盾，因为这里的正体积积分和无边界Stokes用到了紧致性。
\end{proof}
第9章已算出$m>1$时$H^2_{\mathrm{dR}}(S^{2m})=0$，所以这些球面都不能承载Kahler结构；这排除的是Kahler度量，并不排除可能的非Kahler复结构。第17.4节将据此区分拓扑障碍与复可积性问题。


\Needspace{16\baselineskip}
\section{可积性、球面复结构问题的准确边界}
\subsection{Nijenhuis张量与必要条件}
对几乎复结构定义
\begin{equation}\label{eq:nijenhuis}
 N_J(X,Y)=[JX,JY]-J[JX,Y]-J[X,JY]-[X,Y].
\end{equation}
用$[fX,Y]=f[X,Y]-Y(f)X$及$J^2=-I$展开，四个导数项两两相消，故$N_J(fX,Y)=fN_J(X,Y)$；反对称性给另一变量也线性。因此它是张量，尽管定义中出现非张量的李括号。

\begin{proposition}\label{prop:nijenhuis-involutive}
$N_J=0$当且仅当$T^{0,1}$的光滑截面在括号下封闭。复图册诱导的$J$必满足这两个条件。
\end{proposition}
\begin{proof}
复线性延拓$N_J$。对$U,V\in T^{0,1}$，代入$JU=-iU,JV=-iV$得
\[
 N_J(U,V)=-2[U,V]+2iJ[U,V]=-4\pi^{1,0}[U,V].
\]
对两个$(1,0)$向量为共轭式；对一个$(1,0)$和一个$(0,1)$向量直接代入则恒为零。按直和分解展开任意两向量，遂得等价。若有全纯坐标，$T^{0,1}$由可交换的$\partial_{\bar j}$生成，任意系数的括号仍在其张成空间，故满足条件。
\end{proof}

\textbf{Newlander--Nirenberg定理（外引，未在本书证明）。}在光滑几乎复流形上，$N_J=0$保证每点附近存在全纯适配坐标，因而$J$来自复流形结构。上面的命题只证明代数判据及必要方向；充分性涉及复偏微分方程，不能将第5章\emph{实}Frobenius定理直接套在$T^{0,1}$上。它与共轭丛互补但不是一个实切分布，实Frobenius给实积分片的结论没有直接提供所需的复坐标。

\begin{figure}[htbp]
\centering\DGDiagram{complex-conditions}

\caption{本章的条件关系。$N_J=0$到复图册的方向使用外引Newlander--Nirenberg；Kahler与$\nabla J=0$的等价只在已可积的Hermite背景下证明。闭基本形式本身不足以替代可积性。}
\label{fig:complex-conditions}
\end{figure}

\subsection{一个可直接计算的非可积结构}
在$\mathbb C^2$的底层实空间上，定义$T^{0,1}$由
\[
 \bar Z_1=\partial_{\bar z^1},\qquad
 \bar Z_2=\partial_{\bar z^2}+\bar z^1\partial_{z^1}
\]
张成，并令$T^{1,0}$为其共轭。比较$\partial_{z^2}$和$\partial_{\bar z^2}$系数，再比较其余两个系数，即可看出两丛交为零，故处处直和为整个复化切丛。令$J$在两丛上分别为$-i,+i$，它与复共轭交换，因此来自一个光滑实张量，且$J^2=-I$。

然而$[\bar Z_1,\bar Z_2]=\partial_{z^1}=Z_1$，于是$N_J(\bar Z_1,\bar Z_2)=-4Z_1\ne0$。不存在诱导这一$J$的复图册；不是说底层$\mathbb R^4$本身不能具有\emph{其他}可积结构。它当然还有标准复结构。

\subsection{球面问题的三个不同层次}
首先，奇数实维空间不能有$J^2=-I$：实线性空间若成为复向量空间，实维必须为偶数。$S^2\simeq\mathbb{CP}^1$则有本章显式构造的复结构及Kahler度量。

其次，对$m>1$的$S^{2m}$，本章已由上同调和Stokes\emph{证明}不存在Kahler结构。至于哪些偶维球面能有几乎复结构，经典结论是正维球面只有$S^2$和$S^6$，依赖Borel--Serre等拓扑方法；本书未证明该分类，也未构造八元数给出的$S^6$标准几乎复张量。即使某个特定张量不可积，也不能推出底流形上所有几乎复张量都不可积。

\textbf{研究动态与核查范围（2026年11月2日）。}早期讲义把$S^6$是否存在复结构列为Hopf问题，不能将这句话当作不随时间变化的定理。Philip Engel的\textit{Complex structures on $S^6$}（arXiv:2609.38442v1，2026年9月29日提交）给出一种构造的自足论证；Jeff Viaclovsky的\textit{A two-parameter family of complex structures on $S^6$}（arXiv:2609.33785v1，9月27日提交）陈述两参数构造。本轮核对了两篇官方摘要，并读取Engel作者版引言与证明路线；本章将这些记录为\emph{外部研究结果与待专题核验的构造}，没有逐条验证其代数几何、单连通性、整同调与光滑球面识别证明，亦不从预印本页面推断期刊审稿状态。

因此这里不作“截至今日仍未解决”的断言，也不把这些新构造写成本章已证定理。它们与“不存在Kahler结构”没有逻辑冲突：若采用所述复结构，它必属于非Kahler范围。一般Kodaira嵌入、Kahler恒等式与Hodge型分解、Calabi--Yau定理、球面几乎复分类及上述新构造的完整证明，都尚未在本章建立。


\Needspace{18\baselineskip}
\section{复几何习题及详细解答}
17.1、17.2、17.3、17.4、17.8分别改述或扩展Moroianu第1.5节题2/3、第5.5节题1/2、第4.3节题4、第7.4节题2、第1.5节题4；印刷/PDF页分别为10、31、26、41、11。17.5--17.7为原创补充。所有尺度与Chern符号按本书约定，周书页码和梁书对应题尚待核对。

\Needspace{12\baselineskip}
\subsection*{习题17.1：复行列式、定向与反全纯映射}
改述Moroianu第1.5节题2、3并扩展。证明可逆复矩阵的实行列式为复行列式模平方，并说明几乎复流形为何可定向。复共轭在$\mathbb C^m$上何时保持实定向？是否保持复结构？
\begin{proof}[解答]
把复矩阵$C=A+iB$看成实矩阵$C_{\mathbb R}=\left(\begin{smallmatrix}A&-B\\B&A\end{smallmatrix}\right)$，复化后以$z=x+iy,\bar z=x-iy$换基，便把它化为$\operatorname{diag}(C,\bar C)$。因此$\det_{\mathbb R}C_{\mathbb R}=\det C\det\bar C=|\det C|^2>0$。源题使用的实分块次序反转了$B$，行列式结论相同。

局部复基确定实基$(v_1,Jv_1,\ldots,v_m,Jv_m)$，重叠时换基正行列式使定向一致。复共轭在每对实坐标上为$(x^j,y^j)\mapsto(x^j,-y^j)$，所以总实行列式为$(-1)^m$；$m$偶数时保持定向。但它总满足$\dd f\,J=-J\,\dd f$，在$m\geq1$时不全纯。保持实定向比保持复结构弱得多。
\end{proof}

\Needspace{12\baselineskip}
\subsection*{习题17.2：复曲线与共形变化}
取Moroianu第5.5节题1的\emph{已可积}情形，并结合题2扩展。证明复维一Hermite度量必为Kahler；再证明连通复维$m\geq2$的Kahler流形中，$e^{2u}g$仍为Kahler当且仅当$u$为常数。为什么不能将第一句中的可积性凭空删去？
\begin{proof}[解答]
基本形式为实$(1,1)$形式，而实维二时三形式为零，所以$\dd\omega=0$。若只给二维几乎复结构，则还需可积性论证；原讲义已使用Newlander--Nirenberg。本题把起点限定为复曲线，不把该外引隐藏在维数计数里。

共形后$\widetilde\omega=e^{2u}\omega$，故闭性等价于$\dd u\wedge\omega=0$。将其再楔$\omega^{m-2}$得$\dd u\wedge\omega^{m-1}=0$。在正交基中，$u^j\wedge\omega^{m-1}$和$v^j\wedge\omega^{m-1}$分别是仅缺$v^j$和$u^j$的基顶减一次形式，系数均为非零的$(m-1)!$，因此该映射可逆，$\dd u=0$。连通性给$u$为常数；反向常倍显然保持闭性。
\end{proof}

\Needspace{12\baselineskip}
\subsection*{习题17.3：张量积、对偶和Chern联络}
改述Moroianu第4.3节题4并扩展。给定两个带Hermite度量的全纯丛，写出直和、张量积及对偶的Chern联络与曲率，并说明为何是Chern联络而非仅仅某个联络。
\begin{proof}[解答]
直和的联络矩阵为$A_E\oplus A_F$，曲率为$F_E\oplus F_F$。张量积中
\[
 \nabla(s\otimes t)=\nabla^Es\otimes t+s\otimes\nabla^Ft,
 \qquad F_{E\otimes F}=F_E\otimes I+I\otimes F_F.
\]
其矩阵$A_E\otimes I+I\otimes A_F$的平方中，两组交叉项因一形式反交换而抵消，故得到曲率式。对偶从$\dd(\lambda(s))=(\nabla\lambda)(s)+\lambda(\nabla s)$得到$A_{E^*}=-A_E^{\mathsf T}$，曲率$F_{E^*}=-F_E^{\mathsf T}$；转置时矩阵值一形式的楔积次序也须反转，不能只机械地转置$A\wedge A$。

这些联络都保持相应的直和、张量积或对偶内积，且全纯标架中所有势为$(1,0)$型，所以$(0,1)$部分等于自然的$\bar\partial$。定理\ref{thm:chern-connection}的唯一性遂说明它们恰为相应Chern联络。特别地$\det E$的曲率为$\operatorname{tr}F_E$，典范丛的曲率为$-\operatorname{tr}F_{T^{1,0}}$。
\end{proof}

\Needspace{12\baselineskip}
\subsection*{习题17.4：射影直线的真实尺度}
改述Moroianu第7.4节题2，改用显式坐标而不调用常曲率空间整体分类。证明本书$(\mathbb{CP}^1,g_{\mathrm{FS}})$等距于半径$1/\sqrt2$的圆球，求其面积和Ricci形式，并与第16章Hopf底球的单位尺度比较。
\begin{proof}[解答]
第16章映射
\[
 q(w)=\frac{(2\operatorname{Re}w,2\operatorname{Im}w,1-|w|^2)}{1+|w|^2}
\]
全局延拓为到单位球面的微分同胚。直接求导并将映射乘$1/\sqrt2$，得到
\begin{align*}
 q^*g_{S^2}&=\frac{4|\dd w|^2}{(1+|w|^2)^2},\\
 (q/\sqrt2)^*g_{S^2(1/\sqrt2)}&=\frac{2|\dd w|^2}{(1+|w|^2)^2}=g_{\mathrm{FS}}.
\end{align*}
这里第二个拉回中的映射以半径$1/\sqrt2$球面为目标，故给全局等距。面积为$2\pi$，$K=2$，所以$\rho=2\omega_{\mathrm{FS}}$。

Hopf丛前章取单位底球，面积为$4\pi$且$F=i\mathrm{vol}_{S^2}/2$。本章$\omega_{\mathrm{FS}}=\mathrm{vol}_{S^2}/2$，故仍有$F=i\omega_{\mathrm{FS}}$，归一化周期$\int F/(2\pi i)=1$保持不变。度量面积和曲率周期是不同对象，不能同时把所有常数都设成1。
\end{proof}

\Needspace{12\baselineskip}
\subsection*{习题17.5：几乎复情形的平方零失效}
原创补充。对第17.4节的非可积例子，令$\theta^1=\dd z^1-\bar z^1\dd\bar z^2$、$\theta^2=\dd z^2$。验证其为$(1,0)$余标架，并按型别投影定义$\bar\partial_J$，计算$\bar\partial_J^2z^1$。
\begin{proof}[解答]
两形式在$\bar Z_1,\bar Z_2$上均为零，且$\theta^a(Z_b)=\delta^a_b$，所以为所述余标架。其共轭为$\bar\theta^1=\dd\bar z^1-z^1\dd z^2$和$\bar\theta^2=\dd\bar z^2$。于是
\[
 \dd z^1=\theta^1+\bar z^1\bar\theta^2,\qquad
 \bar\partial_Jz^1=\bar z^1\bar\theta^2.
\]
又$\dd\bar z^1=\bar\theta^1+z^1\theta^2$、$\dd\bar\theta^2=0$，取$(0,2)$分量得到
\[
 \bar\partial_J^2z^1=\bar\theta^1\wedge\bar\theta^2\ne0.
\]
同时$\dd\theta^1=-\dd\bar z^1\wedge\dd\bar z^2$有非零$(0,2)$部分，与括号$[\bar Z_1,\bar Z_2]=Z_1$相对应。这个反例说明投影定义的算子未必构成复形。
\end{proof}

\Needspace{12\baselineskip}
\subsection*{习题17.6：重言线丛的全纯标架与酉标架}
原创补充。在$\mathbb{CP}^1$的$w$卡上，重言线丛取全纯标架$e_0=(1,w)$，内积由$\mathbb C^2$限制。求Chern联络、曲率、酉标架下的局部势及$iF/(2\pi)$的积分。
\begin{proof}[解答]
全纯标架的度量系数$H=1+|w|^2$，故
\[
 A_{\mathrm{hol}}=\frac{\bar w\,\dd w}{1+|w|^2},\qquad
 F=-\frac{\dd w\wedge\dd\bar w}{(1+|w|^2)^2}=i\omega_{\mathrm{FS}}.
\]
酉标架为$u_0=e_0/\sqrt H$，换标架$G=H^{-1/2}$不是全纯，但可使用一般联络换元律，得到
\[
 A_u=A_{\mathrm{hol}}-\tfrac12\dd\log H
 =\frac{\bar w\dd w-w\dd\bar w}{2(1+|w|^2)}.
\]
这正是第16章$A_0$，为纯虚形式。全纯标架中的$A_{\mathrm{hol}}$则不必纯虚，因为它不是酉标架。最后$iF=-\omega_{\mathrm{FS}}$，所以$\int iF/(2\pi)=-1$，而$\int F/(2\pi i)=1$。两个归一化相差负号；本题直接算形式，不诉诸一般Chern类理论。
\end{proof}

\Needspace{12\baselineskip}
\subsection*{习题17.7：闭形式的幂与球面障碍}
原创补充。紧致无边界实$2m$维流形上，若实二形式$\omega$闭且非退化，证明可用$\omega^m$定向，并且$[\omega]^k\ne0$。将结论用于$S^6$，解释为何与复结构构造不是同一个问题。
\begin{proof}[解答]
非退化斜型在每点有线性标准基，$\omega=\sum_j u^j\wedge v^j$，故$\omega^m$无零，给出光滑定向。按此定向$\omega^m$处处正，紧致性给有限正积分。若某$\omega^k=\dd\eta$，则$\omega^m=\dd(\eta\wedge\omega^{m-k})$，Stokes使积分零，矛盾。证明甚至不要求已有$J$。

$S^6$的$H^2_{\mathrm{dR}}$为零，故不存在这样的闭非退化二形式，尤其不存在Kahler度量。一个复结构只先规定可积$J$；平均法仍能给Hermite度量，但其基本形式不必闭。因此排除Kahler结构不能用于否定第17.4节所述外部复结构构造。
\end{proof}

\Needspace{12\baselineskip}
\subsection*{习题17.8：Hopf型商与不闭的基本形式}
改述Moroianu第1.5节题4并扩展。设$\alpha>1$实数，$m\geq2$。证明$\alpha^\mathbb Z$在$\mathbb C^m\setminus\{0\}$上的商为紧复流形，底层微分同胚于$S^1\times S^{2m-1}$。证明$|z|^{-2}g_{\mathrm{Euclid}}$下降为Hermite度量，但该度量不是Kahler。是否已证明商上\emph{任何}Hermite度量都非Kahler？
\begin{proof}[解答]
非零点的径向模长在$\alpha^k$作用下改变为$\alpha^k|z|$，故作用自由。任一紧集的模长有正下界和有限上界，仅有限个$k$可能使其平移与自身相交；每点还可选足够薄的径向开环域，使不同平移不交。这样的开集投影给复坐标，换元为常数复线性伸缩。利用极坐标
\[
 z\longmapsto(\log|z|,z/|z|)\in\mathbb R\times S^{2m-1},
\]
群作用成为第一因子的平移$\log\alpha$。其商为$(\mathbb R/(\log\alpha)\mathbb Z)\times S^{2m-1}$，直接给Hausdorff、第二可数、紧致及所述微分同胚，不预用一般适当群作用商定理。

伸缩拉回Euclid度量乘$\alpha^{2k}$，而$|z|^{-2}$乘$\alpha^{-2k}$，故乘积不变，可下降。基本形式为$\widetilde\omega=|z|^{-2}\omega_0$，于是
\[
 \dd\widetilde\omega=-2|z|^{-2}\dd\log|z|\wedge\omega_0.
\]
$\dd\log|z|$处处不为零，复维至少二时前面的楔积单射性说明右式不为零。所以该\emph{特定}度量非Kahler。若要证明商上完全不存在Kahler度量，还需进一步的拓扑或复几何障碍；本题没有计算其第二上同调，不能把度量反例自动升级为流形的不存在定理。
\end{proof}

